% Current v17 support consequences. All graph counts are necessary candidates.
\subsection{全五頂点への伝播と，残る十一型}
\label{sf17:section}

\begin{corollary}[三角形の必要性]
\label{sf17:triangle-required}
正の局所整数 charts，marked puncture の整数 clutching，対向 $D$ の相殺，
および node 非負性を満たす，五次の標準 $K_5$ branch graph を考える．
高分母 support $H=\{ij:q_{ij}>1\}$ が非空なら，$H$ は三角形を含む．
従って $H$ が triangle-free なら，全ての分母は一である．
\end{corollary}

\begin{proof}
系~\ref{fo13:remaining-support} により，辺を持つ成分は一つで，
四または五頂点を使い，閉路を含む．$H$ に三角形がないと仮定する．
五角形を含む場合，その対角線はどれも三角形を作るので，
$H=C_5$ となり，系~\ref{cf13:degree-two-rigidity} に反する．
従って四角形を含む．その二本の対角線も加えられない．
第五頂点から四角形への高辺が二本以上あれば，その行先は隣接できないので，
ちょうど二本で，行先は対向する二頂点である．
可能な形は $C_4$，葉辺付き $C_4$，$K_{2,3}$ に尽くされるが，
それぞれ系~\ref{cf13:degree-two-rigidity}，
定理~\ref{cp14:no-pendant-square}，定理~\ref{kb15:no-K23} に反する．
\end{proof}

\begin{corollary}[残る十一型と，全五頂点の連結性]
\label{sf17:remaining-support}
前系の仮定の下で，高分母 support $H$ が非空なら，
$H$ は五頂点全てを使う連結 graph である．
その第一 Betti 数は二以上であり，少なくとも六辺を持つ．
最大次数は三以上で，孤立頂点は存在しない．
同値な形で，一つの core の四分母が全て一なら，全ての core の全分母が一となる．

ここまでの graph 障害を通過する型は，標識を忘れると十一，
$K_5$ の標識を保つと三百五十六である．
これらは必要条件を通過した候補の個数であり，
整数 charts の同時接着や幾何的実現の存在を主張しない．
\end{corollary}

\begin{proof}
辺を持つ成分は一つであることは系~\ref{fo13:remaining-support} による．
系~\ref{kg17:low-core-propagation} により，非空の場合に孤立頂点はない．
従って全五頂点が一つの連結成分に属する．

閉路を含むので第一 Betti 数は一以上である．
それが一なら，前系によって唯一の閉路は三角形となる．
三角形に長さ二の道を付ける場合，または同じ頂点に二本の葉辺を付ける場合は，
三角形の二頂点が二価のままなので
定理~\ref{tn14:no-degree-two-triangle} に反する．
残る，異なる二頂点に葉辺を一本ずつ付けた形も
定理~\ref{tg16:no-two-leaf-triangle} により排除された．
従って第一 Betti 数は二以上であり，辺数は少なくとも六である．
五頂点で六辺以上を持つので最大次数は三以上である．

個数は $K_5$ の全 $2^{10}=1024$ 部分 graph に対する完全な有限検算で得られる．
既存の十二型・三百六十一標識から，
定理~\ref{kg17:no-K4} により $G_{4,3}=K_4$ の五標識を除くと，
十一型・三百五十六標識となる．
$K_4$ の使用頂点は五頂点から四つを選ぶので，除く標識数は $\binom54=5$ とも確認できる．
全 graph の列挙，同型判定，旧識別子の保持，各代表の自己同型群の検算は
\nolinkurl{verify_support_agenda_v17.py} と
\nolinkurl{support_agenda_verification_v17.json} に保存した．
この列挙は分母や行列 entry の有界探索ではない．
\end{proof}

局所的には，二価 core の二つの高分母近傍の間が低辺なら，
系~\ref{kb15:negative-closes-triangle} によりその core は正の枝に属する．
さらに任意の高辺と共通低 puncture に対して，
通常接着なら補題~\ref{cl16:ordinary-bound} の reverse bound $|u|\le4$ が成立する．
これら二つの局所的な制約には，対向 $D$ の相殺は不要である．

ここで除いたのは，support 全体が四使用頂点上の $K_4$ である場合である．
$K_4$ を部分 graph として含み，第五頂点にも高辺を持つ型に対して，
定理~\ref{kg17:no-K4} をそのまま適用することはできない．
例えば $G_{5,12}$ は $K_4$ に一本の葉辺を加えた型であり，
その第五 core は一高分母となるため，引き続き候補に含める．
